\chapter*{Lời Tựa}
\addcontentsline{toc}{chapter}{Lời Tựa}
\markboth{Lời Tựa}{Lời Tựa}

\begin{truyen}{Lời Tựa}{Foreword}
\chuhoa{Q}{uyển sách này được hình thành không chỉ từ nỗi lo tài chính của hiện tại, mà còn từ một ký ức rất dài của thiếu thốn. Có một đứa trẻ bảy tuổi mất cha, không có nhà cửa ổn định, theo gia đình lưu lạc nhiều nơi, từng đi chăn vịt, từng vào lớp một khi đã bảy tuổi.}
\textit{(This book was formed not only from today’s financial worries, but from a long memory of deprivation. There was once a seven-year-old child who lost his father, had no stable home, wandered with the family from place to place, herded ducks, and entered first grade already at the age of seven.)}

Đứa trẻ ấy nay đã ba mươi chín tuổi, trở thành một người thầy đứng lớp, có vợ, có con, có đồng lương và có những buổi dạy thêm để gồng gánh gia đình. Nhưng người từng đi qua đói nghèo hiểu rõ rằng chỉ cần buông tay trong chi tiêu, cái khổ cũ có thể quay lại dưới một hình thức khác.
\textit{(That child is now thirty-nine years old, a teacher with a classroom, a wife, children, a salary, and extra lessons taken on to carry the family. Yet one who has lived through poverty knows that if spending discipline is loosened, old hardship may return in another form.)}

Đây không phải là một giáo trình tài chính. Nó là tập hợp những câu chuyện nhỏ, những nỗi day dứt rất thật, những điển tích xưa được kéo lại gần đời sống hôm nay để nhắc mỗi người rằng tiết kiệm không phải thói quen của người yếu, mà là bản lĩnh của người biết lo xa. Với người viết, đây còn là cuốn sách để tự nhắc mình đừng phụ tuổi thơ đã chịu quá nhiều cơ cực; với học sinh, đây là cuốn sách mong các em học sớm một điều mà nhiều người lớn hiểu quá muộn: biết giữ đồng tiền là biết giữ nhân cách và tương lai của mình.
\textit{(This is not a formal finance manual. It is a collection of small stories, real anxieties, and classical anecdotes brought close to modern life to remind us that saving is not a habit of the weak, but a strength of those who think ahead. For the writer, it is also a book reminding him not to betray the hardship of his own childhood; for students, it is a hope that they may learn early what many adults understand too late: to keep money wisely is also to protect character and future.)}

Nếu người đọc tìm thấy chính mình trong những trang này, xin đừng ngại. Bởi ai cũng từng có lúc buông tay với tiền bạc, chỉ khác nhau ở chỗ ta có chịu dừng lại để sửa hay không.
\textit{(If readers find themselves in these pages, they should not feel ashamed. Everyone has at some point been careless with money; the difference lies in whether we stop and correct ourselves.)}
\end{truyen}

\begin{baihoc}
Biết tự răn mình trước khi trách hoàn cảnh là bước đầu tiên để sống vững hơn. Người từng khổ mà vẫn còn biết sợ sự phung phí sẽ có cơ hội dựng lại mái nhà của mình chắc hơn.
\textit{(Knowing how to correct oneself before blaming circumstances is the first step toward a steadier life. Those who have suffered yet still fear wastefulness have a better chance to rebuild a stronger home.)}
\end{baihoc}

\begin{ghinhoanh}
\textbf{Short English to remember:}

Self-correction comes before stability.

Discipline is a quiet form of courage.
\end{ghinhoanh}

\begin{flushright}
\textbf{Biên soạn cho một người thầy muốn tự răn mình và dạy học sinh biết sống chắt chiu}\\
\textit{Tháng 3 năm 2026}
\end{flushright}

\ngancach